
\section{Extracting the constituents of verbal constructions in UD}\label{UD}

We assume a certain familiarity with the UD annotation scheme in this paper and simply recall that the syntactic structure of all sentences is encoded as a dependency tree. Here we focus on verbal constructions, that is, a verb and all the constituents that depend on it. We consider all lexical verbs and have decided not to restrict our study to main verbs (see (\ref{ex:notmain})) for at least two reasons: first, including all lexical verbs can double or triple the number of constructions, which is not negligible when the corpus is small; second, in some sentences, the most interesting verb in terms of construction may be subordinated because the main verb is merely a modal or an auxiliary.

All dependents of a verb (in the tree) cannot be considered constituents (of the verbal construction we are interested in). The following are excluded in our data extraction: punctuation signs (attached by the \texttt{punct} relation in UD), discourse markers (\texttt{discourse} relation; see \textit{alas} in (\ref{ex:aux})), juxtaposed clauses or parentheses (\texttt{parataxis} relation), conjuncts in a coordination (\texttt{conj} relation; (\ref{ex:coord})), and coordinating conjunctions (\texttt{cc} relation; see \textit{but} in (\ref{ex:aux})). The \texttt{vocative} relation, which concerns addresses to a participant (\textit{Guys, take it easy!}), is excluded but dislocated arguments (\textit{These guys, I don't trust them.}) are kept in, especially because in some languages the boundary between governed units and dislocated units is difficult to draw. The \texttt{aux} relation (see \textit{do} in (\ref{ex:aux})), is excluded, considering that auxiliaries are part of the verb form; compounds (see \textit{out} in (\ref{ex:out})) are excluded as well. Markers of subordination such as subordinating conjunctions and adpositions (\texttt{mark} and \texttt{case} relations; see \textit{when} in (\ref{ex:notmain})), which are analysed as dependents in UD but are generally considered as heads of their constructions, are also discarded.\footnote{In her study, \citet{tanaka2021menzerath} did not exclude any dependent, except the punctuation.}

Examples of English verbal constructions are given in (\ref{ex:1}) for illustration; the verb is underlined and the constituents are shown between square brackets; filtered elements are given in italics.

\pex\label{ex:1}
\a\label{ex:aux}\textit{But} [I] , \textit{alas} , \textit{do} [not] \underline{know} [how to see sheep] [through the walls of boxes]. \corpus{en\_littleprince}
\a\label{ex:notmain}This means that \textit{when} [you] \underline{move} [a series or category field] [to the filter area and back], previously hidden items are again hidden. \corpus{en\_lines}
\a\label{ex:coord}[You] \underline{tell} [me] [what they look like] \textit{and I'll tell you what they are}. \corpus{en\_childes}
\a\label{ex:out}[One minister] [reportedly] \underline{handed} \textit{out} [100 dollar `gifts'] [to journalists attending a press conference for Allawi], [a practice that brings back bad memories to many Iraqis]. \corpus{en\_ewt}
\xe

Another possibility taken into account is that a constituent may be split into two contiguous parts, for example in cases of extraction or extraposition. These cases are treated as two separate constituents. For instance, in (\ref{ex:what}), the relative pronoun \textit{what} is extracted and separated from its governor \textit{do}, and in (\ref{ex:extra}), the relative clause \textit{that go well with the Santorum cocktail} is extraposed and separated from its governor \textit{dishes}.

\pex
\a\label{ex:what}[What] \textit{do} [you] \underline{like} [to do] ? \corpus{en\_ewt} 
\a\label{ex:extra}[What culinary dishes] \textit{would} [you] \underline{recommend} [that go well with the Santorum cocktail]? \corpus{en\_gum}
\xe

For each language L and each $n \geq$ 1, we extracted all the verbal constructions with $n$ constituents. We then calculated the size of each constituent, excluding punctuation. In other words, we kept all relations except \texttt{punct}; if a constituent contains a coordination or a discourse marker, there is no reason to remove them. 

We call MAL$_n$(L) the mean size of the constituents in all verbal constructions with $n$ constituents. Note that we compute MAL$_n$(L) for a given language L only when we have at least 100 constructions with $n$ constituents. We return to this point in Section~\ref{metrics}. We define LMAL$_n$(L) and RMAL$_n$(L) in the same way. For LMAL, all constructions with $n$ preverbal constituents (i.e.\ on the left) were selected, regardless of the other side of the verb. Likewise for RMAL, with all constructions having $n$ postverbal constituents (i.e.\ on the right).

